\documentclass[12pt,a4paper]{article}

% ============================================================
% NEJM Review: Physiologic Pacing in Heart Failure
% 整理：謝慕揚 MD, PhD, FESC
% ============================================================

% Drake_preamble.tex - LaTeX Preamble Template
% 作者: 謝慕揚, MD, PhD, FESC
% 
% 使用說明:
% 1. 表格請使用 longtable，不要有垂直分隔線 |
% 2. 表格內換行使用 \newline，不要使用 \\
% 3. 藥物名稱、檢驗、檢查請維持英文
% 4. Reference 格式請使用 \href 加上驗證過的 DOI

% Including essential packages for Chinese typesetting
\usepackage{xeCJK}
\usepackage{fontspec}
% Configuring Chinese font with Noto Serif CJK TC, fallback to SimSun
\IfFontExistsTF{Noto Serif CJK TC}{
  \setCJKmainfont{Noto Serif CJK TC}
  \setCJKsansfont{Noto Serif CJK TC}
  \setCJKmonofont{Noto Serif CJK TC}
}{\setCJKmainfont{SimSun}
  \setCJKsansfont{SimSun}
  \setCJKmonofont{SimSun}
}

% Including other necessary packages
\usepackage{geometry}
\usepackage{titlesec}
\usepackage{enumitem}
\usepackage{xcolor}
\usepackage{colortbl}
\usepackage{tcolorbox}
\usepackage{booktabs}
\usepackage{longtable}
\usepackage{array}
\usepackage{graphicx}
\usepackage{tikz}
\usepackage{fancyhdr}
\usepackage{multicol}
\usepackage{datetime2}
\usepackage{hyperref}
\usepackage{mdframed}
\usepackage{amsmath}
\usepackage{amssymb}
\usepackage{multirow}

\hypersetup{
    colorlinks=true,
    linkcolor=emergency_blue,
    citecolor=emergency_blue,
    urlcolor=emergency_blue,
    pdfborder={0 0 0}
}

% Defining page geometry
\geometry{
  top=2.5cm,
  bottom=2.5cm,
  left=2cm,
  right=2cm,
  headheight=25pt
}

% Adjusting topmargin to compensate for increased headheight
\addtolength{\topmargin}{-10pt}

% Defining colors
\definecolor{emergency_red}{RGB}{186,24,27}
\definecolor{emergency_blue}{RGB}{0,114,188}
\definecolor{emergency_gray}{RGB}{120,120,120}
\definecolor{highlight_yellow}{RGB}{255,250,205}

% Setting title formats
\titleformat{\section}[block]{\Large\bfseries\color{emergency_red}}{\thesection}{10pt}{}
\titleformat{\subsection}[block]{\large\bfseries\color{emergency_blue}}{\thesubsection}{8pt}{}
\titleformat{\subsubsection}[block]{\normalsize\bfseries\color{emergency_gray}}{\thesubsubsection}{6pt}{}

% Defining custom environment for important notes
\newmdenv[
    linecolor=emergency_red,
    backgroundcolor=highlight_yellow,
    linewidth=2pt,
    topline=true,
    bottomline=true,
    leftline=true,
    rightline=true,
    innertopmargin=10pt,
    innerbottommargin=10pt,
    innerleftmargin=10pt,
    innerrightmargin=10pt
]{importantbox}

% Defining note command with tcolorbox
\newcommand{\note}[1]{
    \par
    \noindent
    \begin{tcolorbox}[
        colback=emergency_blue!5!white,
        colframe=emergency_blue,
        title=\textbf{重點提醒},
        fonttitle=\bfseries,
        boxrule=0.5pt,
        arc=3pt,
        left=6pt,
        right=6pt,
        top=6pt,
        bottom=6pt
    ]
    #1
    \end{tcolorbox}
}

% Key findings box
\newcommand{\keyfinding}[1]{
    \par
    \noindent
    \begin{tcolorbox}[
        colback=yellow!10!white,
        colframe=orange!75!black,
        title=\textbf{核心發現摘要},
        fonttitle=\bfseries,
        boxrule=0.5pt,
        arc=3pt
    ]
    #1
    \end{tcolorbox}
}

% Defining custom date format for ISO 8601
\DTMnewdatestyle{iso}{
  \renewcommand{\DTMdisplaydate}[4]{\number##1-\number##2-\number##3}
  \renewcommand{\DTMDisplaydate}[4]{\number##1-\number##2-\number##3}
}
\DTMsetdatestyle{iso}
\newcommand{\mydate}{\DTMtoday}


% Custom header
\pagestyle{fancy}
\fancyhf{}
\fancyhead[L]{\textcolor{emergency_red}{NEJM 教學講義}}
\fancyhead[R]{\textcolor{emergency_blue}{Physiologic Pacing in Heart Failure}}
\fancyfoot[C]{\textcolor{emergency_gray}{\thepage}}
\renewcommand{\headrulewidth}{1pt}
\renewcommand{\headrule}{\hbox to\headwidth{\color{emergency_red}\leaders\hrule height \headrulewidth\hfill}}

\begin{document}

%============================================================
% Title Page
%============================================================

\begin{center}
{\LARGE\bfseries\textcolor{emergency_red}{The New England Journal of Medicine}}\\[0.5cm]
{\Large\textbf{Physiologic Pacing in Heart Failure}}\\[0.3cm]
{\large 心臟衰竭之生理性心律調節治療}\\[0.5cm]
{\normalsize \href{https://doi.org/10.1056/NEJMra2415650}{N Engl J Med 2026;394:367-81. DOI: 10.1056/NEJMra2415650}}\\[0.3cm]
{\normalsize 整理：謝慕揚 MD, PhD, FESC \quad 日期：\mydate}
\end{center}

\vspace{0.5cm}

%============================================================
% Key Findings Box
%============================================================

\keyfinding{
\textbf{重大發現：}
\begin{itemize}[leftmargin=*,nosep]
    \item Cardiac resynchronization therapy (CRT) 現稱為 cardiac physiologic pacing，適用於 HFrEF 合併傳導系統疾病患者
    \item Biventricular pacing 在超過 10,000 名心衰患者的臨床試驗中證實可改善 LVEF、減少二尖瓣逆流與心室容積
    \item 最大受益族群：LBBB + QRS $\geq$150 msec；QRS $\leq$130 msec 時效益有限
    \item Conduction system pacing (His bundle/left bundle branch pacing) 為新興替代策略，回顧性研究顯示可能優於 biventricular pacing
\end{itemize}

\textbf{臨床意義：}傳導系統起搏正在改變心衰合併傳導異常的治療模式，多項大型隨機試驗進行中。
}

\tableofcontents
\newpage

%============================================================
\section{文獻概述}
%============================================================

本篇為 2026 年 1 月發表於 NEJM 的 Review Article，系統性回顧心臟衰竭之生理性心律調節治療 (physiologic pacing)，涵蓋傳統 biventricular pacing 與新興 conduction system pacing 策略。

\subsection{文獻資訊}

\begin{itemize}
    \item \textbf{標題：}Physiologic Pacing in Heart Failure
    \item \textbf{作者：}Chelu MG, Poole JE, Ellenbogen KA
    \item \textbf{機構：}Baylor College of Medicine, University of Washington, Virginia Commonwealth University
    \item \textbf{期刊：}N Engl J Med 2026;394:367-81
    \item \textbf{DOI：}\href{https://doi.org/10.1056/NEJMra2415650}{10.1056/NEJMra2415650}
\end{itemize}

%============================================================
\section{基本概念}
%============================================================

\subsection{名詞定義}

Cardiac resynchronization therapy (CRT) 現於學會指引中更名為 \textbf{cardiac physiologic pacing}，以更準確描述其生理機制。

\subsection{適應症}

心臟生理性起搏適用於：
\begin{enumerate}
    \item LVEF $\leq$ 50\% 之心衰患者
    \item 預期或已有高比例心室起搏負荷 (high ventricular pacing burden)
    \item 寬 QRS complex (傳導系統疾病)
\end{enumerate}

\subsection{技術分類}

\begin{longtable}{p{4cm}p{10cm}}
\toprule
\textbf{技術類型} & \textbf{特點} \\
\midrule
\endfirsthead
\toprule
\textbf{技術類型} & \textbf{特點} \\
\midrule
\endhead
\bottomrule
\endfoot

\textbf{Biventricular Pacing} &
• 傳統 CRT 方式\newline
• 右心室導線 + 冠狀竇分支導線\newline
• 透過融合 RV 與 LV 電波前線達到再同步化\newline
• 豐富的 RCT 證據支持 \\
\midrule

\textbf{His Bundle Pacing (HBP)} &
• 最生理性的起搏方式\newline
• 保留左右心室正常活化順序\newline
• His bundle 結構小，定位困難\newline
• 閾值可能升高或變化 \\
\midrule

\textbf{Left Bundle Branch Pacing (LBBP)} &
• 導線穿過心室中隔至左束支\newline
• 較大的目標區域，成功率較高\newline
• 閾值較低且穩定\newline
• 可矯正 His bundle 以下的 LBBB \\
\midrule

\textbf{HOT-CRT} &
• His bundle-optimized CRT\newline
• His bundle pacing + 冠狀竇導線 \\
\midrule

\textbf{LOT-CRT} &
• Left bundle branch-optimized CRT\newline
• Left bundle pacing + 冠狀竇導線 \\
\bottomrule
\end{longtable}

%============================================================
\section{Biventricular Pacing 臨床證據}
%============================================================

\subsection{早期試驗}

早期小型試驗 (MUSTIC-SR, MIRACLE, PATH-CHF 等) 證實 biventricular pacing 可改善：
\begin{itemize}
    \item LVEF 與左室容積 (reverse remodeling)
    \item 二尖瓣逆流程度
    \item 生活品質與運動耐受度
    \item NYHA functional class
\end{itemize}

\subsection{大型隨機試驗結果}

\begin{longtable}{p{3cm}p{2cm}p{4cm}p{5cm}}
\toprule
\textbf{試驗} & \textbf{人數} & \textbf{主要終點} & \textbf{結果} \\
\midrule
\endfirsthead
\toprule
\textbf{試驗} & \textbf{人數} & \textbf{主要終點} & \textbf{結果} \\
\midrule
\endhead
\bottomrule
\endfoot

\textbf{COMPANION} & 1,520 & 全因死亡或任何原因住院 & CRT ± ICD 優於單純藥物治療\newline QRS $\geq$150 msec 獲益最大 \\
\midrule

\textbf{CARE-HF} & 813 & 全因死亡或心血管住院 & CRT 顯著改善複合終點\newline 長期追蹤降低 sudden cardiac death\newline QRS $\geq$160 msec 效果最佳 \\
\midrule

\textbf{MADIT-CRT} & 1,820 & 全因死亡或心衰事件 & CRT-D 優於 ICD 單獨\newline QRS $\geq$150 msec + LBBB 獲益最大\newline Non-LBBB 無顯著效益 \\
\midrule

\textbf{RAFT} & 1,798 & 全因死亡或心衰住院 & CRT 降低心衰事件\newline LBBB + QRS $\geq$150 msec 獲益明顯 \\
\midrule

\textbf{REVERSE} & 610 & 心衰臨床複合評分 & CRT 改善臨床評分\newline QRS $\leq$130 msec 效益有限 \\
\midrule

\textbf{BLOCK HF} & 691 & 全因死亡、心衰急診或 LVESI 增加 $\geq$15\% & Biventricular pacing 優於 RV pacing\newline 適用於 AV block + LVEF $\leq$50\% \\
\bottomrule
\end{longtable}

\subsection{窄 QRS 患者試驗}

\note{
\textbf{重要發現：}四項針對 QRS $<$120-130 msec 患者的試驗 (NARROW-CRT, LESSER-EARTH, EchoCRT, RethinQ) 均未顯示 CRT 效益。
\begin{itemize}
    \item 窄 QRS 患者不建議植入 CRT
    \item 心臟超音波機械性不同步無法預測 CRT 反應
\end{itemize}
}

\subsection{Biventricular Pacing 無反應者}

約 \textbf{1/3} 患者對 biventricular pacing 無反應，可能原因包括：

\begin{itemize}
    \item Non-LBBB QRS morphology
    \item QRS 120-130 msec (邊緣寬度)
    \item 缺乏適當冠狀竇靶血管
    \item 高起搏閾值導致失去 capture
    \item 末期心衰
    \item 競爭性心律 (頻繁 PVC 或傳導性 AF) 導致起搏比例 <95\%
    \item 側壁心肌瘢痕
\end{itemize}

\textbf{反應較佳族群：}
\begin{itemize}
    \item 非缺血性心肌病優於缺血性
    \item 女性在較短 QRS 即可獲益
    \item 身材較矮者
\end{itemize}

%============================================================
\section{Conduction System Pacing}
%============================================================

\subsection{傳導系統解剖}

\begin{itemize}
    \item \textbf{His bundle：}長度約 2.6 mm，寬度 3.7 mm，厚度 1.4 mm，穿過膜性中隔，被中央纖維體包圍
    \item \textbf{Left bundle branch：}較寬的心內膜下結構，寬度可達 14 mm，分為前分支、後分支，常有中隔分支
\end{itemize}

\subsection{His Bundle Pacing}

\begin{itemize}
    \item 最生理性的起搏方式，保留正常左右心室活化
    \item 可用於 AV nodal 與 intra-Hisian 傳導疾病
    \item 可矯正部分 LBBB（病變位於 His-Purkinje 進入左束支處）
    \item \textbf{挑戰：}目標區小、定位困難、閾值可能升高、microdislodgement 風險
\end{itemize}

\subsection{Left Bundle Branch Pacing}

\begin{itemize}
    \item 導線穿過心室中隔至左束支或其分支
    \item \textbf{優勢：}
    \begin{itemize}
        \item 成功率高（MELOS 研究：bradycardia 92.4\%，心衰 82.2\%）
        \item 閾值較低且穩定
        \item 可矯正 His bundle 以下的 LBBB
    \end{itemize}
    \item \textbf{併發症：}急性穿孔至左心室 (3.7\%)，但可重新定位，長期無明顯不良影響
\end{itemize}

\subsection{Conduction System Pacing 臨床試驗}

\begin{longtable}{p{3cm}p{2cm}p{3.5cm}p{5.5cm}}
\toprule
\textbf{試驗} & \textbf{人數} & \textbf{比較} & \textbf{結果} \\
\midrule
\endfirsthead
\toprule
\textbf{試驗} & \textbf{人數} & \textbf{比較} & \textbf{結果} \\
\midrule
\endhead
\bottomrule
\endfoot

\textbf{His-SYNC} & 41 & His CRT vs BiV CRT & 高 crossover 率 (48\%)\newline 無顯著差異 \\
\midrule

\textbf{His-Alternative} & 50 & His CRT vs BiV CRT & ITT 無差異\newline Per-protocol: His pacing LVEF 改善較多 \\
\midrule

\textbf{LBBP-RESYNC} & 40 & LBBP vs BiV pacing & LBBP: LVEF 改善較多、LVESV 較小\newline NT-proBNP 較低 \\
\midrule

\textbf{LEVEL-AT} & 70 & CSP vs BiV pacing & LV activation time 無顯著差異\newline 6 個月結果相似 \\
\midrule

\textbf{ALTERNATIVE-AF} & 50 & His pacing vs BiV pacing (AVNA 後) & His pacing: LVEF 改善較多 \\
\midrule

\textbf{HOPE-HF} & 167 & His pacing vs no pacing & Peak VO$_2$ 與 LVEF 無差異\newline 生活品質改善 \\
\bottomrule
\end{longtable}

\subsection{觀察性研究}

\begin{longtable}{p{3cm}p{2cm}p{9cm}}
\toprule
\textbf{研究} & \textbf{人數} & \textbf{主要發現} \\
\midrule
\endfirsthead
\toprule
\textbf{研究} & \textbf{人數} & \textbf{主要發現} \\
\midrule
\endhead
\bottomrule
\endfoot

\textbf{I-CLAS (LVEF $\leq$35\%)} & 1,778 & LBBAP vs BiV pacing: 全因死亡 + 心衰住院複合終點較低、LVEF 改善較多、QRS 較窄 \\
\midrule

\textbf{I-CLAS (LVEF 36-50\%)} & 1,004 & CSP vs BiV pacing: 複合終點較低 \\
\midrule

\textbf{I-CLAS (心律不整)} & 1,414 & LBBAP vs BiV pacing: 持續性 VT/VF 較少、新發 AF 較少 \\
\midrule

\textbf{Zhu et al.} & 259 & LBBP 優於 BiV pacing 和 LV septal pacing\newline LVSP 死亡率較高 \\
\bottomrule
\end{longtable}

\note{
\textbf{關鍵發現：}Meta-analysis (4 RCTs + 17 觀察性研究) 顯示 conduction system pacing 相較於 biventricular pacing 可能降低全因死亡率。多項大型 RCTs 進行中將釐清相對效益。
}

%============================================================
\section{進行中的大型隨機試驗}
%============================================================

\begin{longtable}{p{3cm}p{2cm}p{4cm}p{5cm}}
\toprule
\textbf{試驗} & \textbf{人數} & \textbf{比較} & \textbf{主要終點} \\
\midrule
\endfirsthead
\toprule
\textbf{試驗} & \textbf{人數} & \textbf{比較} & \textbf{主要終點} \\
\midrule
\endhead
\bottomrule
\endfoot

\textbf{LEAP} & 470 & LVSP vs RVP & 死亡、心衰住院、LVEF 下降 $>$10\% \\
\midrule

\textbf{OptimPacing} & 683 & LBBP vs RVP & 死亡、心衰住院、PICM \\
\midrule

\textbf{LEFT HF} & 1,280 & LBBP vs RVP & CV 死亡、心衰事件、LVESV 增加 15\% \\
\midrule

\textbf{PROTECT-HF} & 2,600 & His/LBBP vs RVP & 死亡、心衰事件 \\
\midrule

\textbf{Left vs. Left RCT} & 2,136 & His/LBBP vs BiV pacing & 死亡、心衰住院 \\
\bottomrule
\end{longtable}

%============================================================
\section{臨床實務考量}
%============================================================

\subsection{心臟衰竭與電生理聯合門診}

多專科團隊合作模式有助於：
\begin{itemize}
    \item 優化指引導向藥物治療 (GDMT)
    \item 選擇適當的生理性起搏策略
    \item 改善病患預後、增加參與度、減少住院率與死亡率
\end{itemize}

\subsection{MRI 相容性}

部分 conduction system pacing 導線和系統已取得 MRI conditional 標籤，進行中的臨床試驗將提供更多安全性數據。

\subsection{導線拔除}

\begin{itemize}
    \item Lumenless leads 拔除成功率 100\%（1\% 有殘留遠端碎片）
    \item 輕微併發症 2.1\%
    \item 平均植入時間 22±26 個月，相對較短
    \item Stylet-driven leads 拔除經驗有限
\end{itemize}

%============================================================
\section{未來展望}
%============================================================

\begin{enumerate}
    \item \textbf{技術選擇：}若 conduction system pacing 與 biventricular pacing 效果相當，其他因素將影響選擇：
    \begin{itemize}
        \item 系統較簡單
        \item 手術時間較短
        \item 輻射與顯影劑劑量較低
        \item 電池壽命較長
        \item 成本較低（使用一般 pacemaker/ICD generator，僅需一條導線）
    \end{itemize}

    \item \textbf{待解決問題：}
    \begin{itemize}
        \item His bundle pacing 是否會被 LBBP 完全取代？
        \item Conduction system pacing 對 RBBB 或 nonspecific IVCD 是否有效？
        \item LBBB + 心衰是否以 conduction system pacing 為首選，biventricular pacing 為備案？
        \item 部分患者可能受益於同時 conduction system pacing + CS lead (HOT-CRT/LOT-CRT)
    \end{itemize}
\end{enumerate}

%============================================================
\section{重點整理}
%============================================================

\begin{tcolorbox}[
    colback=emergency_red!5!white,
    colframe=emergency_red,
    title=\textbf{Key Points: Physiologic Pacing in Heart Failure},
    fonttitle=\bfseries
]
\begin{enumerate}
    \item \textbf{適應症：}心衰 + LVEF $\leq$50\% + 高心室起搏負荷或寬 QRS

    \item \textbf{Biventricular pacing 證據：}>10,000 名患者的 RCTs 證實改善 LVEF、減少 MR 與 LV 容積、改善生活品質與運動能力、減少心衰住院、改善存活率

    \item \textbf{最大受益族群：}LBBB + QRS $\geq$150 msec

    \item \textbf{無效族群：}QRS $\leq$130 msec

    \item \textbf{Conduction system pacing：}回顧性研究顯示可能優於 biventricular pacing，多項 RCTs 進行中

    \item \textbf{LBBP vs His bundle pacing：}LBBP 成功率較高、閾值較穩定，正逐漸成為主流
\end{enumerate}
\end{tcolorbox}

%============================================================
\section{參考文獻}
%============================================================

\begin{enumerate}
    \item Chelu MG, Poole JE, Ellenbogen KA. Physiologic Pacing in Heart Failure. \href{https://doi.org/10.1056/NEJMra2415650}{\textit{N Engl J Med}. 2026;394:367-81.}

    \item Chung MK, Patton KK, Lau C-P, et al. 2023 HRS/APHRS/LAHRS guideline on cardiac physiologic pacing for the avoidance and mitigation of heart failure. \href{https://doi.org/10.1016/j.hrthm.2023.03.1538}{\textit{Heart Rhythm}. 2023;20(9):e17-e91.}

    \item Bristow MR, Saxon LA, Boehmer J, et al. Cardiac-resynchronization therapy with or without an implantable defibrillator in advanced chronic heart failure. \href{https://doi.org/10.1056/NEJMoa032423}{\textit{N Engl J Med}. 2004;350:2140-50.}

    \item Cleland JGF, Daubert J-C, Erdmann E, et al. The effect of cardiac resynchronization on morbidity and mortality in heart failure. \href{https://doi.org/10.1056/NEJMoa050496}{\textit{N Engl J Med}. 2005;352:1539-49.}

    \item Moss AJ, Hall WJ, Cannom DS, et al. Cardiac-resynchronization therapy for the prevention of heart-failure events. \href{https://doi.org/10.1056/NEJMoa0906431}{\textit{N Engl J Med}. 2009;361:1329-38.}

    \item Wang Y, Zhu H, Hou X, et al. Randomized trial of left bundle branch vs biventricular pacing for cardiac resynchronization therapy. \href{https://doi.org/10.1016/j.jacc.2022.07.019}{\textit{J Am Coll Cardiol}. 2022;80:1205-16.}

    \item Vijayaraman P, Sharma PS, Cano Ó, et al. Comparison of left bundle branch area pacing and biventricular pacing in candidates for resynchronization therapy. \href{https://doi.org/10.1016/j.jacc.2023.05.006}{\textit{J Am Coll Cardiol}. 2023;82:228-41.}
\end{enumerate}

%============================================================
\section{縮寫對照表}
%============================================================

\begin{longtable}{p{4cm}p{10cm}}
\toprule
\textbf{縮寫} & \textbf{全名} \\
\midrule
\endfirsthead
\toprule
\textbf{縮寫} & \textbf{全名} \\
\midrule
\endhead
\bottomrule
\endfoot

AF & Atrial Fibrillation (心房顫動) \\
AV & Atrioventricular (房室) \\
BiV & Biventricular (雙心室) \\
CRT & Cardiac Resynchronization Therapy (心臟再同步治療) \\
CS & Coronary Sinus (冠狀竇) \\
CSP & Conduction System Pacing (傳導系統起搏) \\
HBP & His Bundle Pacing (希氏束起搏) \\
HFrEF & Heart Failure with Reduced Ejection Fraction (射出分率降低型心衰) \\
HOT-CRT & His Bundle-Optimized CRT \\
ICD & Implantable Cardioverter-Defibrillator (植入式心臟去顫器) \\
IVCD & Intraventricular Conduction Delay (心室內傳導延遲) \\
LBBB & Left Bundle Branch Block (左束支傳導阻滯) \\
LBBP & Left Bundle Branch Pacing (左束支起搏) \\
LOT-CRT & Left Bundle Branch-Optimized CRT \\
LV & Left Ventricle (左心室) \\
LVEF & Left Ventricular Ejection Fraction (左室射出分率) \\
LVESV & Left Ventricular End-Systolic Volume (左室收縮末期容積) \\
NYHA & New York Heart Association \\
PICM & Pacing-Induced Cardiomyopathy (起搏誘發心肌病變) \\
RBBB & Right Bundle Branch Block (右束支傳導阻滯) \\
RV & Right Ventricle (右心室) \\
RVP & Right Ventricular Pacing (右心室起搏) \\
VT/VF & Ventricular Tachycardia/Ventricular Fibrillation \\
\bottomrule
\end{longtable}

\end{document}
